% (c) Meta Platforms, Inc. and affiliates. Confidential and proprietary.
\documentclass[12pt]{article}

\usepackage{amsmath, amssymb, amsthm}
\usepackage{graphicx}
\usepackage{booktabs}
\usepackage{natbib}
\usepackage[margin=1in]{geometry}
\usepackage{hyperref}

\title{The Crucial Role of Multicalibration in Model-based Prevalence Measurement}
\author{
  Fridolin Linder\thanks{Meta Platforms, Inc.} \and
  Thomas Leeper\thanks{Meta Platforms, Inc.}
  % TODO: confirm author order and affiliations
}
\date{\today}

\begin{document}
\maketitle

\begin{abstract}
% TODO: write abstract
\end{abstract}

%%%%%%%%%%%%%%%%%%%%%%%%%%%%%%%%%%%%%%%%%%%%%%%%%%%%%%%%%%%%%%%%%%%%%%%%%%%%%%%
\section{Introduction}
%%%%%%%%%%%%%%%%%%%%%%%%%%%%%%%%%%%%%%%%%%%%%%%%%%%%%%%%%%%%%%%%%%%%%%%%%%%%%%%

Measuring the prevalence of a category within a target population is a task common across many scientific applications. This method includes use of diagnostic tests to measure disease prevalence, measuring prevalence of phenomena on online platforms, and identifying the frequency with which documents in a corpus address a specific topic. Prevalence estimation - or “quantification” in computer science terminology - typically involves using an imperfect measurement device to classify examples into categories.

Recently, large language models (LLMs) have dramatically expanded the scope and scale of this enterprise. Researchers now routinely deploy LLMs as zero-shot measurement devices, without any task-specific training data, to estimate the prevalence of phenomena that previously required expensive manual annotation. Applications span the social sciences, where LLMs code democracy indicators across countries \citep{weidmann2025democracy}, classify protest events in news corpora \citep{overos2024protest}, and categorize open-ended survey responses at near-human accuracy \citep{mellon2024survey, gilardi2023chatgpt}; clinical medicine, where they extract diagnostic attributes from pathology reports \citep{sushil2024clinical} and identify goals-of-care discussions in clinical notes \citep{sibley2025goc}; and the digital humanities, where they annotate art forms in auction records \citep{tojima2025art} and convert qualitative text into quantitative variables across multiple languages \citep{karjus2025mixed}. LLMs are also benign advertised for content moderation (https://openai.com/index/using-gpt-4-for-content-moderation/) and so-called LLM “judges” are used to assess the quality and safety of AI systems ([TODO: Pick citation]). Yet the ease of deployment belies a fundamental methodological gap: virtually none of these studies address what happens to prevalence estimates when the target population differs from the population on which the device's error properties were assessed.

Quantification is a well-studied problem, and several methods exist to adjust for imperfect devices (Gonzalez et al.). However, as we demonstrate, all of these devices can fail to recover correct prevalence estimates in common settings. Nearly all existing methods attempt to learn or calibrate the imperfect device by learning the device’s error properties on a source population of examples and apply that calibrated device to the target population(s). However, in situations where the target population differs in composition from the source population this calibration procedure does not protect against heavily biased prevalence estimates.

We propose multicalibration as the necessary and sufficient condition for accurate, out-of-domain prevalence estimation. Building on recent theoretical advances in algorithmic fairness and domain adaptation (Kim et al., 2022), we leverage the fact that a multicalibrated classifier achieves 'Universal Adaptability' to solve the prevalence estimation problem. We show that this theoretical guarantee allows for unbiased estimates on any downstream target population without the need for target-specific re-weighting or retraining. This robustness is critical for nearly all practical applications of quantification, including applying a model to previously unseen datasets, tracking changes in prevalence over time, or analyzing subsets of the training population. In contrast, we show that simple calibration is insufficient for unbiased estimation except in trivial and unrealistic settings.

We review existing approaches to the quantification problem, connect the quantification task to the broader theory of domain adaptation via multicalibration, demonstrate the need for multicalibration to achieve unbiased estimates of prevalence in out-of-domain settings, and apply those insights to a few real-world datasets.


%%%%%%%%%%%%%%%%%%%%%%%%%%%%%%%%%%%%%%%%%%%%%%%%%%%%%%%%%%%%%%%%%%%%%%%%%%%%%%%
\section{AI Systems as Measurement Devices / Alternative: Imperfect Measurement Devices}
%%%%%%%%%%%%%%%%%%%%%%%%%%%%%%%%%%%%%%%%%%%%%%%%%%%%%%%%%%%%%%%%%%%%%%%%%%%%%%%

[TODO: Explicit Review, maybe split in Application papers and Methodology papers]

%%%%%%%%%%%%%%%%%%%%%%%%%%%%%%%%%%%%%%%%%%%%%%%%%%%%%%%%%%%%%%%%%%%%%%%%%%%%%%%
\section{Calibrated Measurement Devices are Biased under Distribution Shift}
%%%%%%%%%%%%%%%%%%%%%%%%%%%%%%%%%%%%%%%%%%%%%%%%%%%%%%%%%%%%%%%%%%%%%%%%%%%%%%%

Assume a binary problem with outcome y in {0, 1} and a set of relevant features X. The researcher relies on an imperfect device, h(X), such as an AI system, diagnostic test, a human annotator, or a machine learning classifier. The task of "quantification" is to estimate the true prevalence, P(y=1), in a target population using only the device’s outputs. It is well understood that naively averaging the model’s predictions will lead to biased estimates of prevalence (Gonzales ? or find other citations). Several methods are commonly used to adjust or calibrate a model to produce unbiased estimates. In this section, we will explore some of these methods and demonstrate how they fail if the target population has a different distribution of features X (covariate or mix-shift).

[TODO: State assumptions. Causality X->Y, no unobserved confounders, stable conditional probabilities P(Y=1|X), etc.]

%----------------------------------------------------------------------
\subsection{Existing Calibration Approaches}
%----------------------------------------------------------------------

The most intuitive approach, "Classify & Count," simply tallies the device's positive predictions. Such predictions are obtained by thresholding if the model produces a continuous score. For many AI systems that produce natural language output a binary (e.g. yes/no) answer is the native output. However, this method is fundamentally flawed; even highly accurate devices yield biased estimates whenever the false positive and false negative rates are non-zero and uncorrected. To address this, the quantification literature has historically relied on error correction methods, such as the "Adjusted Count" (AC) (In epidemiology, this correction is widely known as the Rogan-Gladen estimator). These methods use a classifier's fixed error rates (tpr and fpr) estimated from training data to invert the confusion matrix and recover the true prevalence (Gonzales et al. 2017)
\citep{gonzalez2017review}.%

[TODO: [discuss / mention other non-classification methods (King; Resnick; etc.)]]

Adjusted Count is a \emph{calibration} technique. If a device is perfectly
calibrated (i.e., its predicted probabilities match empirical frequencies),
averaging its predictions yields an unbiased prevalence estimate.
Reformulating Adjusted Count, we can treat the device's classifications as
real numbers rather than discrete predicted values $\{0,1\}$, which
correspond to the prevalence within each predicted class. Positively
classified examples are scored by the TPR and negatively classified
examples are scored by $1 - \text{FPR}$. These scores satisfy the
definition of calibration that $\mathbb{E}[y \mid \hat{y}] = y$. The
Adjusted Count (Rogan--Gladen) estimate is equivalent to the mean of these
scores.
% TODO: verify this reformulation is correct---see Fridolin's margin comment
% questioning whether the algebra works out.

A prediction $\hat{y}$ is \emph{perfectly calibrated} if and only if
\[
  P(y = 1 \mid \hat{y} = p) = p,
\]
where $p$ is the predicted probability. Intuitively, among all instances
receiving score $\hat{y} = 0.8$, we expect 80\% to have label~$1$.

% TODO: note that recent review (Grimmer, Roberts, Stewart) hardly mentions
% calibration at all \citep{grimmer2022text}.

\paragraph{Failure under mix shift.}
Corrections such as AC rely on a strict and often unrealistic assumption:
while the class prevalence $P(y)$ may change, the distribution of features
within each class, $P(X \mid y)$, remains constant between training and
target populations. This ``Feature Stability'' assumption fails in most
real-world out-of-distribution settings. For example, if a device is
applied to a specific demographic subgroup or a future time period, the
features characterizing the positive class often shift, altering the
device's error rates and rendering the standard correction invalid
\citep{wu2024stable}.

This can be demonstrated for AC. Consider a population consisting of two
groups, indexed by~$i$, of sizes $\{k, 1-k\}$. Each group is defined by a
prevalence~$p_i$ and error rates $\text{FPR}_i$ and $\text{FNR}_i$. The
apparent overall FPR and FNR of the population is a weighted average of the
error rates for the two groups. In the presence of a ``mix shift'' that
changes the relative sizes of these groups to $\{k^*, 1-k^*\}$, the
learned TPR and TNR will no longer be an appropriately weighted average of
the target population group-level error rates. Estimated prevalence---corrected
by an inappropriate global confusion matrix---will be biased. Despite
achieving ``calibration'' in the original population, this learned
calibration does not always hold.
% TODO: analogy here to conditioning on post-treatment variables
% (comment from Thomas Leeper)

% TODO: something about why we should care about bias (versus variance)

While theoretically sound, reliance on global calibration introduces a
subtle but critical failure mode. A device can be \emph{globally}
calibrated on a training set---accurate on average---while being grossly
\emph{miscalibrated} on specific subgroups, even within the training
distribution. When researchers apply such a model to a new domain---whether
a specific demographic subset, a new geographic region, or a future time
period---the global calibration curve often fails to hold. This failure
occurs because most real-world distribution shifts are effectively changes
in the composition of the population. If the device is miscalibrated on the
underlying subgroups that make up the population, any shift in their
proportions will bias the aggregate estimate. This phenomenon, which we
term \emph{mis-multicalibration}, violates the Calibration Stability
assumption locally and negatively impacts generalization. Consequently,
neither standard error-correction methods (which assume stable features)
nor standard calibration methods (which assume stable global calibration)
can yield unbiased estimates under shift.

%----------------------------------------------------------------------
\subsection{Simulation Illustration}
%----------------------------------------------------------------------

To illustrate our theoretical results, we use a simulation with the
following structure. We generate synthetic data with a binary covariate
$X \in \{0, 1\}$ and binary outcome $Y \in \{0, 1\}$, where
$P(Y=1 \mid X=1) = 0.85$ and $P(Y=1 \mid X=0) = 0.15$. A classifier
produces probability estimates~$\hat{p}$ that are systematically
miscalibrated: underestimating by 10\% when $X=0$ and overestimating by
10\% when $X=1$. All calibration parameters are learned on a training
distribution with $P(X=0) = 0.5$. We then evaluate estimation methods on
test distributions where $P(X=0)$ ranges from 0.01 to 0.99, representing
covariate shift. The marginal distribution $P(X)$ changes while the
conditional $P(Y \mid X)$ remains fixed. The simulation is repeated
50~times.

% TODO: insert Figure 1 - 4-panel figure with panels a, b, c, d
\begin{figure}[htbp]
  \centering
  % \includegraphics[width=\textwidth]{figures/simulation_bias.pdf}
  \caption{Bias of the global prevalence estimate under covariate shift and
    three adjustment methods. (a)~No adjustment, averaging raw model
    predictions. (b)~Averaging binarized model predictions from a
    (globally) calibrated threshold. (c)~Averaging predictions calibrated
    on a calibration sample. (d)~Adjusted Count (Rogan--Gladen adjustment)
    with TPR and FPR estimated from a calibration sample.}
  \label{fig:simulation-bias}
\end{figure}

Figure~\ref{fig:simulation-bias} displays the results for (a)~the raw
simulated model prediction, (b)~a classify-and-count procedure with a
calibrated threshold, (c)~probability calibrated model scores, and (d)~the
adjusted-count method (Rogan--Gladen). The $y$-axis shows the bias in
global prevalence estimates under varying covariate shift, which is
displayed on the $x$-axis. With zero distribution shift (the center of the
figure), the raw score average is biased by about 7\%. The
classify-and-count, calibration, and the Rogan--Gladen estimators show no
bias when there is no distribution shift, but their prevalence estimates
are biased if the distribution shifts. The threshold-based estimators (CC,
and RG) are much more prone to extreme bias---in this sample up to
250\%---as the covariate shift becomes more extreme. The calibrated scores
are less susceptible to such bias amplification.
% TODO: make all y-axes the same scale (comment from Thomas Leeper)
% TODO: add a classify-and-count example that might look the same
% (comment from Thomas Leeper)

%----------------------------------------------------------------------
\subsection{Multicalibration}
%----------------------------------------------------------------------

To guarantee accurate quantification in out-of-domain settings, a device
must satisfy multicalibration: it must be calibrated not just on average,
but simultaneously across all practically relevant subpopulations. By
ensuring the device is reliable on these ``atomic'' components,
multicalibration provides robustness not just for subsetting, but for
generalizing to any future distribution composed of these groups
\citep{gopalan2022omnipredictors}.

Multicalibration was initially introduced as a fairness criterion
\citep{hebertjohnson2018multicalibration, blasiok2023loss} and a growing
body of research has shown its relevance to several aspects of model
performance and robustness. Formally, let $X$ denote features, $Y \in
\{0,1\}$ an outcome, and $f(X) \in [0,1]$ a predictor interpreted as a
probability. Given a collection $\mathcal{G}$ of subgroups $G \subseteq
\mathcal{X}$ (for example, groups defined by demographic attributes or by
arbitrary computable functions of~$X$), $f$ is said to be
$\alpha$-multicalibrated with respect to $\mathcal{G}$ if, for every group
$G \in \mathcal{G}$ and every prediction value~$v$ used by the model:
\[
  \left| \mathbb{E}[Y \mid f(X)=v,\, X \in G] - v \right| \le \alpha,
\]
whenever the conditioning event has sufficient probability mass. Within
each subgroup and at each score level, the average observed outcome closely
matches the predicted probability. This notion strengthens classical
calibration, which requires the condition only over the entire population,
by enforcing reliability across a rich family of subpopulations,
potentially exponential in size.

% TODO: fill in how this applies to measurement and illustrate with simulation

% TODO: insert multicalibration simulation figure
% \begin{figure}[htbp]
%   \centering
%   \includegraphics[width=0.6\textwidth]{figures/multical_simulation.pdf}
%   \caption{Multicalibration simulation results.}
%   \label{fig:multical-simulation}
% \end{figure}

\paragraph{Empirically achievable multicalibration.}
% TODO: discuss feature selection / heterogeneity / data size availability
% TODO: results across varying levels of model quality

Mix shift can resemble ``concept drift'' where concept drift has in fact
not occurred. Multicalibration mitigates the risk of mix-shift-induced
bias, but does not address true concept drift (change in meaning of
classes).

Recent theoretical advances \citep{wu2024stable} propose relaxing the
rigid requirement of Feature Stability in favor of ``Calibration
Stability.'' Under this framework, accurate measurement does not require
the features~$X$ to behave identically across populations; it only requires
the relationship between the device's score~$s$ and the outcome~$y$---the
calibration curve $P(y \mid s)$---to remain stable.

% TODO: results under imperfect multicalibration; sensitivity analysis:
% shifting distribution, correlation with outcome, numbers/selection of features

% TODO: discuss true concept drift---how do we set this aside?

%%%%%%%%%%%%%%%%%%%%%%%%%%%%%%%%%%%%%%%%%%%%%%%%%%%%%%%%%%%%%%%%%%%%%%%%%%%%%%%
\section{Empirical Applications}
%%%%%%%%%%%%%%%%%%%%%%%%%%%%%%%%%%%%%%%%%%%%%%%%%%%%%%%%%%%%%%%%%%%%%%%%%%%%%%%

% TODO: select and implement 1--2 empirical examples.
% Candidates:
%   - Disease prevalence measurement with noisy tests (Rogan-Gladen etc.)
%     - Diabetes screening in youth (Vangeepuram et al., 2021)
%     - Survey-based self-reported HIV status (Metz et al., 2024)
%     - Cause of death misclassification (Edwards et al., 2025)
%   - Phenomena on social media (Weidmann & Roel, 2023)
%   - Document topics
%   - Vangeepuram et al. has publicly available data and specifically
%     assesses screening test quality by subgroup.
% TODO: decide on datasets

%%%%%%%%%%%%%%%%%%%%%%%%%%%%%%%%%%%%%%%%%%%%%%%%%%%%%%%%%%%%%%%%%%%%%%%%%%%%%%%
\section{Conclusion}
%%%%%%%%%%%%%%%%%%%%%%%%%%%%%%%%%%%%%%%%%%%%%%%%%%%%%%%%%%%%%%%%%%%%%%%%%%%%%%%

% TODO: write conclusion
% Key point: Multicalibration has mostly been applied in the context of
% fairness, but has much broader implications/applications.

%%%%%%%%%%%%%%%%%%%%%%%%%%%%%%%%%%%%%%%%%%%%%%%%%%%%%%%%%%%%%%%%%%%%%%%%%%%%%%%
% References
%%%%%%%%%%%%%%%%%%%%%%%%%%%%%%%%%%%%%%%%%%%%%%%%%%%%%%%%%%%%%%%%%%%%%%%%%%%%%%%
\bibliographystyle{apalike}
\bibliography{references}

%%%%%%%%%%%%%%%%%%%%%%%%%%%%%%%%%%%%%%%%%%%%%%%%%%%%%%%%%%%%%%%%%%%%%%%%%%%%%%%
\appendix
\section{Simulation Study Setup}
\label{app:simulation}
%%%%%%%%%%%%%%%%%%%%%%%%%%%%%%%%%%%%%%%%%%%%%%%%%%%%%%%%%%%%%%%%%%%%%%%%%%%%%%%

\subsection{Data Generating Process}

We generate synthetic data with a binary covariate $X \in \{0, 1\}$ and
binary outcome $Y \in \{0, 1\}$. The conditional outcome probabilities
are:
\[
  P(Y=1 \mid X=1) = 0.85, \qquad P(Y=1 \mid X=0) = 0.15.
\]

A simulated classifier produces probability estimates~$\hat{p}$ that are
systematically miscalibrated within each stratum:
\begin{align*}
  \text{For } X=0: \quad \hat{p} &= P(Y=1 \mid X=0) \times 0.9 = 0.135
    \quad \text{(10\% underestimation)}, \\
  \text{For } X=1: \quad \hat{p} &= P(Y=1 \mid X=1) \times 1.1 = 0.935
    \quad \text{(10\% overestimation)}.
\end{align*}
This represents a classifier whose predictions are directionally correct
but exhibit stratum-specific bias---a common pattern in real-world machine
learning systems.

\subsection{Covariate Shift Protocol}

All calibration parameters are learned on a training distribution with
$P(X=0) = 0.5$. We then evaluate prevalence estimation accuracy across
shifted test distributions where $P(X=0)$ varies from 0.01 to 0.99. This
range represents the full spectrum of potential distribution shifts, from
populations dominated by $X=1$ to those dominated by $X=0$.

The key assumption is that of covariate shift: while the marginal
distribution $P(X)$ changes between training and test, the conditional
distribution $P(Y \mid X)$ remains constant.

\subsection{Prevalence Estimation Methods}

We compare five methods for estimating population prevalence
$\pi = P(Y=1)$:

\paragraph{1. Uncalibrated Averaging.}
\[
  \hat{\pi}_{\text{uncal}} = \frac{1}{n}\sum_{i=1}^{n} \hat{p}_i.
\]
The raw average of predicted probabilities without any correction.

\paragraph{2. Classify and Count (CC).}
\[
  \hat{\pi}_{\text{CC}} = \frac{1}{n}\sum_{i=1}^{n}
    \mathbb{1}[\hat{p}_i \geq \tau].
\]
Predictions are binarized at threshold~$\tau$, chosen on calibration data
to match true prevalence, then the positive proportion is computed.

\paragraph{3. Rogan--Gladen Adjustment.}
\[
  \hat{\pi}_{\text{RG}} = \frac{\bar{p} - \text{FPR}}
    {\text{TPR} - \text{FPR}},
\]
where $\text{TPR} = \mathbb{E}[\hat{p} \mid Y=1]$ and
$\text{FPR} = \mathbb{E}[\hat{p} \mid Y=0]$ are estimated from
calibration data.

\paragraph{4. Global Calibration.}
\[
  \hat{\pi}_{\text{cal}} = \alpha \cdot \bar{p}, \qquad \text{where }
  \alpha = \frac{\bar{Y}_{\text{cal}}}{\bar{p}_{\text{cal}}}.
\]
A multiplicative correction factor learned to match expected prevalence on
calibration data.

\paragraph{5. Multicalibration.}
\[
  \hat{\pi}_{\text{MC}} = \frac{1}{n}\sum_{i=1}^{n}
    (\hat{p}_i + \delta_{X_i}), \qquad \text{where }
  \delta_x = \mathbb{E}[Y \mid X=x] - \mathbb{E}[\hat{p} \mid X=x].
\]
Stratum-specific additive corrections that ensure calibration conditional
on~$X$.

\subsection{Simulation Procedure}

For each of $B=50$ bootstrap iterations:
\begin{enumerate}
  \item Generate fresh calibration data ($n=10{,}000$) from the training
    distribution ($P(X=0)=0.5$).
  \item Estimate all calibration parameters from this data.
  \item For each of 20 test distributions with
    $P(X=0) \in [0.01, 0.99]$:
    \begin{enumerate}
      \item Generate test data ($n=10{,}000$).
      \item Apply each estimation method using the learned calibration
        parameters.
      \item Compute percentage bias:
        $\text{Bias}\% = 100 \times
          \frac{\hat{\pi} - \pi_{\text{true}}}{\pi_{\text{true}}}$.
    \end{enumerate}
\end{enumerate}
This bootstrap procedure properly accounts for estimation variance in
calibration parameters.

\subsection{Results}

Multicalibration produces unbiased prevalence estimates across all
distribution shifts (bias $\approx 0\%$ throughout). This robustness
arises because multicalibration ensures
$\mathbb{E}[\hat{p} \mid X] = \mathbb{E}[Y \mid X]$ for each stratum.
Under covariate shift, where $P(Y \mid X)$ is invariant, these
stratum-level corrections remain valid regardless of changes in~$P(X)$.

Global Calibration and Classify-and-Count exhibit similar bias patterns:
zero bias at the training distribution ($P(X=0)=0.5$) but linearly
increasing bias as the distribution shifts. At extreme shifts, bias
reaches approximately $\pm 15\%$.

Uncalibrated estimates show moderate shift sensitivity with bias scaling
approximately linearly with distribution shift magnitude.

Rogan--Gladen adjustment exhibits the most severe bias amplification, with
bias exceeding $\pm 40\%$ at extreme shifts. This instability stems from
two factors: (1)~TPR and FPR parameters estimated on training data become
incorrect under shift, and (2)~the ratio estimator's small denominator
($\text{TPR} - \text{FPR}$) amplifies estimation errors.

\subsection{Implications}

These results demonstrate that for applications where covariate shift is
expected---such as deploying models across populations with different
demographic composition---multicalibration provides the most robust
prevalence estimates. The critical requirement is that calibration
corrections be learned conditional on covariates that may shift at
deployment time. Standard calibration approaches that only ensure marginal
calibration ($\mathbb{E}[\hat{p}] = \mathbb{E}[Y]$) are insufficient
under distribution shift.

\end{document}
